\section{Bài Toán}

\begin{frame}{Tổng quan đề tài}
  \begin{columns}[T,totalwidth=\textwidth]
    \begin{column}{0.46\textwidth}
  \vspace{0.5cm}
      \only<-5>{\begin{itemize}[<+->]
        \item Xây dựng AI chơi Cờ Tướng.
        \item AI được mô hình hóa như một Intelligent Agent.
        \item Quan sát và phân tích trạng thái bàn cờ.
        \item Đưa ra nước đi tốt nhất trong môi trường đối kháng.
        \item Một trong những bài toán kinh điển của AI: Game playing.
      \end{itemize}}
      \only<6->{
        \begin{itemize}
          \item<6-> Cờ Tướng là bài toán:
            \begin{enumerate}
              \item Hai người chơi.
              \item Đối kháng Zero-Sum.
            \end{enumerate}
          \item<7-> Không gian tìm kiếm rất lớn ($\approx 10^{150}$)
          \item<8-> Vì vậy, đòi hỏi:
            \begin{enumerate}
              \item Tìm kiếm thông minh.
              \item Đánh giá chính xác.
              \item Hiệu năng cao.
            \end{enumerate}
        \end{itemize}
      }
    \end{column}
    \begin{column}{0.50\textwidth}
      \begin{figure}
        \centering
        \includegraphics[height=0.64\textheight]{\figpath{fig-00-board-start.png}}
        \caption{Bàn cờ khởi đầu.}
      \end{figure}
    \end{column}
  \end{columns}
\end{frame}

\begin{frame}[t]{Kiến trúc hệ thống}
  \begin{columns}[T,totalwidth=\textwidth]
    \begin{column}{0.4\textwidth}
      \begin{block}{Agent}
        \textbf{Đầu vào}
        \begin{itemize}[<+->]
          \item Trạng thái hiện tại của bàn cờ.
          \item Lượt chơi hiện tại.
          \item Chuỗi các nước đi.
        \end{itemize}
        \onslide<4->{\textbf{Đầu ra}}
        \begin{itemize}[<+->]
          \item Nước đi tốt nhất cho trạng thái bàn cờ hiện tại.
        \end{itemize}
      \end{block}
    \end{column}
    \begin{column}{0.6\textwidth}
      \begin{figure}[htbp]
        \centering
        \definecolor{blueFill}{HTML}{E1EDFF}
        \definecolor{blueStroke}{HTML}{2B70E4}
        \definecolor{greenFill}{HTML}{E6F9EC}
        \definecolor{greenStroke}{HTML}{28A745}
        \definecolor{redFill}{HTML}{FFEBEB}
        \definecolor{redStroke}{HTML}{DC3545}
        \begin{tikzpicture}[
            scale=0.5,
            transform shape,
            % --- Node Styles ---
            blueNode/.style={
              rectangle, 
              draw=blueStroke, 
              fill=blueFill, 
              line width=2pt, 
              rounded corners=4pt,
              align=center,
              inner sep=10pt,
            },
            greenNode/.style={
              rectangle, 
              draw=greenStroke, 
              fill=greenFill, 
              line width=2pt, 
              rounded corners=4pt,
              align=center,
              inner sep=10pt,
            },
            redNode/.style={
              rectangle, 
              draw=redStroke, 
              fill=redFill, 
              line width=2pt, 
              dashed,
              rounded corners=4pt,
              align=center,
              inner sep=10pt,
            },
            % --- Fixed Cylinder Style ---
            dbNode/.style={
              cylinder, 
              draw=greenStroke, 
              fill=greenFill, 
              line width=2pt, 
              shape border rotate=90, 
              align=center, 
              inner sep=10pt,            % Reduced inner padding
              aspect = 0.25
            },
            % --- Arrow Styles ---
            arrow/.style={
              -{Stealth}, 
              draw=black!80
            },
            dashedArrow/.style={
              -{Stealth}, 
              dashed, 
              draw=black!80
            },
            biArrow/.style={
              {Stealth}-{Stealth}, 
              draw=black!80
            },
            ]

            %% --- Nodes Placement ---
            \node (GUI) [blueNode] at (0,-0.92) {GUI};
            \node (ThreadA) [blueNode, right=of GUI, xshift=1.5cm] {Luồng A: Stdin Listener};
            \node (Engine) [greenNode, below=of ThreadA] {Luồng B: UCI Engine};

            % Right Side Node
            \node (State) [redNode, right=of ThreadA, xshift=1.5cm] {Cờ trạng\\thái};

            % Core Engine Stack
            \node (ID) [greenNode, below=of Engine] {Iterative Deepening};
            \node (Algo) [greenNode, below=of ID, xshift=1.5cm] {Negamax Alpha-Beta};

            % Bottom & Left Nodes
            \node (Eval) [greenNode, below=of Algo] {Hàm đánh giá tĩnh};
            \node (Trans) [dbNode, left=2.2cm of Algo] {Transposition\\Table}; % Adjusted spacing slightly for balance

            %% --- Connections & Labels ---
            % GUI -> ThreadA
            \draw [arrow] (GUI) -- node[below, font=\small, align=center] {Lệnh\\qua stdin} (ThreadA);

            % ThreadA -> Engine
            \draw [arrow] (ThreadA) -- node[left, font=\small] {Chuyển lệnh} (Engine);

            % Engine -> GUI (Loop back loop on the left side)
            \draw [arrow] (Engine.west) .. controls +(-3.5,0) and +(0,0) .. node[font=\small, align=center, below, yshift=-0.25cm] {Trả về\\bestmove} (GUI.south);

            % ThreadA -> State
            \draw [arrow] (ThreadA) -- node[below, font=\small, align=center] {Ngắt khẩn\\cấp} (State);

            % Engine -> ID -> Algo
            \draw [arrow] (Engine) -- (ID);
            \draw [arrow] (ID) -- (Algo);

            % State -> Algo (Dashed routing down to the side of Algo)
            \draw [dashedArrow] (State.south) -- node[font=\small, align=left, right, xshift=0.2cm] {Đọc liên tục\\để dừng} (Algo.east);

            % Algo -> Eval
            \draw [arrow] (Algo) -- node[midway, right, font=\small] {} (Eval);

            % Algo <-> Trans (Bidirectional)
            \draw [biArrow] (Algo) -- node[midway, below, font=\small, align=center] {Lưu\\/Đọc thế} (Trans);

          \end{tikzpicture}
          \caption{Sơ đồ kiến trúc của Lingine}
        \end{figure}

    \end{column}
  \end{columns}
\end{frame}

\begin{frame}{Giao thức UCI}
  \begin{figure}
      \centering
        \begin{tikzpicture}[
            scale=1,
            transform shape,
            % Styles for lifelines and boxes
            guiBox/.style={draw=lingineBlue, fill=lingineBlue!10, thick, rounded corners=2pt, minimum width=2.0cm, minimum height=0.6cm, align=center, font=\small},
            engBox/.style={draw=lingineGreen, fill=lingineGreen!10, thick, rounded corners=2pt, minimum width=2.0cm, minimum height=0.6cm, align=center, font=\small},
            lifeline/.style={draw=black!30, dashed, thick},
            % Arrow styles
            toEngine/.style={-{Stealth}, draw=lingineBlue, thick},
            toGui/.style={-{Stealth}, draw=lingineGreen, thick},
            infoLine/.style={-{Stealth}, draw=black!40, dashed, thick},
            % Text style
            lbl/.style={font=\scriptsize\ttfamily, midway, above, yshift=-1pt}
          ]
          
          % Lifeline Headers
          \node[guiBox] (GUI) at (0,0) {Chess GUI};
          \node[engBox] (Engine) at (7.5,0) {Lingine (Engine)};

          \path let \p1 = (GUI) in 
            \pgfextra{\xdef\guix{\x1}}; % Store x-coordinate in a macro
          \path let \p1 = (Engine) in 
            \pgfextra{\xdef\enginex{\x1}}; % Store x-coordinate in a macro
          
          % Lifeline Paths
          \coordinate[below=5.2cm of GUI] (GUIbottom);
          \coordinate[below=5.2cm of Engine] (Enginebottom);
          \draw[lifeline] (GUI) -- (GUIbottom);
          \draw[lifeline] (Engine) -- (Enginebottom);
          
          % Communication Steps
          % Step 1: Handshake
         \onslide<1-> \draw[toEngine] (\guix,-0.7) -- node[lbl] {uci} (\enginex,-0.7);
          \onslide<2-> \draw[toGui] (\enginex,-1.3) -- node[lbl] {uciok} (\guix,-1.3);
          
          % Step 2: Ready Check
          \onslide<3-> \draw[toEngine] (0,-2.0) -- node[lbl] {isready} (\enginex,-2.0);
          \onslide<4-> \draw[toGui] (\enginex,-2.6) -- node[lbl] {readyok} (0,-2.6);
          
          % Step 3: Calculation Flow
       \onslide<5->    \draw[toEngine] (\guix,-3.3) -- node[lbl] {position startpos moves ...} (\enginex,-3.3);
          \onslide<6-> \draw[toEngine] (\guix,-3.9) -- node[lbl] {go depth 10} (\enginex,-3.9);
          
          % Step 4: Search Info & Bestmove
 \onslide<7->          \draw[infoLine] (\enginex,-4.5) -- node[lbl] {info depth 1 ...} (\guix,-4.5);
        \onslide<8->  \draw[toGui] (\enginex,-5.0) -- node[lbl] {bestmove h0g2} (\guix,-5.0);
          
        \end{tikzpicture}
      \caption{Luồng giao tiếp qua giao thức UCI}
    \end{figure}
\end{frame}

% Định nghĩa cột mới tự động canh lề trái, tự giãn độ rộng và thêm \hspace{0pt} để sửa lỗi baseline của \onslide
\newcolumntype{L}[1]{>{\raggedright\arraybackslash\hspace{0pt}}p{#1}}

\begin{frame}{Các trở ngại khi lập trình cờ tướng}
  \begin{columns}[T,totalwidth=\textwidth]
    \begin{column}{0.48\textwidth}
      \begin{tabular}{L{2.2cm} L{3.8cm}}
        \toprule
        Thành phần & Cờ tướng \\[0.25cm]
        \midrule
        % CÁCH SỬA: Đặt \onslide ở đầu dòng và tuyệt đối KHÔNG bao trọn dòng bằng dấu {}
        \onslide<2-> Bàn cờ & $10 \times 9 = 90$ giao điểm rất lớn \\ \addlinespace
        \onslide<3-> Vùng luật & Cung và sông \\ \addlinespace
        \onslide<4-> Cản quân & Chân Mã, mắt Tượng, ngòi Pháo \\ \addlinespace
        \onslide<5-> Hàm đánh giá & Phức tạp, yêu cầu nhiều kiến thức cho tàn cuộc \\ \addlinespace
        \onslide<6-> Bộ luật đặc biệt & Lộ mặt Tướng, chiếu/đuổi dai \\
        \bottomrule
      \end{tabular}
    \end{column}
    \onslide<1->
    \begin{column}{0.50\textwidth}
      \begin{figure}
        \centering
        \includegraphics[height=0.64\textheight]{\figpath{fig-01-blocking.png}}
        \caption{Cản mắt tượng và chân mã.}
      \end{figure}
    \end{column}
  \end{columns}
\end{frame}



